\section{Conclusiones}
\label{sec:conclusiones}

Este Trabajo de Fin de Título ha demostrado la viabilidad técnica de implementar un sistema completo de diagnóstico del vehículo OBD-II. Los resultados obtenidos en el banco de ensayo Arduino, junto con la validación sobre un vehículo real del grupo VAG (Škoda Yeti de 2015), confirman que la plataforma desarrollada es funcional, interoperable con ECUs de producción y extensible en sus tres componentes de forma independiente.

Las conclusiones más relevantes del trabajo se pueden agrupan en torno a los tres componentes del sistema:

\textbf{Sobre la pila de protocolos.} La implementación directa del UDS y OBD-II sobre \texttt{SocketCAN} en Python ha resultado viable. La librería \texttt{can-isotp} proporciona una abstracción suficiente para gestionar la segmentación y el Flow Control de forma transparente, si bien la correcta sincronización de la secuencia multi-frame para el VIN requirió un ajuste cuidadoso de los parámetros de timing en el simulador. La latencia extremo a extremo del sistema (28--62\,ms por consulta) es desprecible pues se trata de una herramienta de diagnóstico para uso personal que no trata factores críticos.

\textbf{Sobre el simulador Arduino.} El simulador de ECU desarrollado se ha revelado como una herramienta de validación de alto valor, que va más allá de la simple emulación de respuestas fijas. El modelo físico del motor, los mecanismos opcionales de variabilidad del bus (ruido en sensores y tramas de difusión) y la inyección de DTCs por interrupción hardware permiten al banco de ensayo la capacidad de simulación de escenarios reales que resulta difícil de reproducir sin acceso continuo a un vehículo.

\textbf{Sobre la arquitectura software.} La decisión de estructurar la herramienta Python con interfaces abstractas e inyección de dependencias ha demostrado alto valor a lo largo del desarrollo: la posibilidad de sustituir el transporte CAN real por un \textit{mock} en cualquier momento redujo significativamente los ciclos de depuración y permitió trabajar con la capa de aplicación de forma independiente del hardware.

Desde el punto de vista del proceso de desarrollo, la metodología basada en prototipos incrementales resultó especialmente adecuada para un proyecto de esta naturaleza, en el que la validez de las hipótesis técnicas no podía garantizarse a priori. Los ciclos cortos de implementación y prueba sobre el banco de ensayo permitieron detectar y corregir problemas de sincronización ISO-TP en etapas tempranas, antes de que estos afectaran al desarrollo de las capas superiores.

La validación sobre vehículo real se llevó a cabo finalmente sobre un Škoda Yeti de 2015 (no sobre el SEAT Ibiza 6J planteado inicialmente, por problemas de acceso ya comentados previamente). Como limitación cabe señalar, por tanto, que la validación se realizó sobre un único vehículo y a través del conector OBD-II, lo que da acceso al bus de diagnóstico.


\section{Líneas de trabajo futuro}
\label{sec:trabajo_futuro}

Los resultados obtenidos y las limitaciones identificadas abren varias líneas de continuación para este trabajo:

\textbf{Validación sobre flota y acceso al bus de tracción.} La validación se ha realizado sobre un único vehículo (Škoda Yeti 2015). Una continuación natural sería ampliarla a una flota de varios modelos de distintos fabricantes y calibrar el modelo físico del simulador con datos reales de cada plataforma.

\textbf{Soporte de servicios UDS adicionales.} El sistema implementa ya los servicios UDS DiagnosticSessionControl (\texttt{0x10}) y ReadDataByIdentifier (\texttt{0x22}). Una extensión es incorporar otros servicios de ISO 14229-1, como la escritura de datos por identificador (\texttt{0x2E}), el control de rutinas (\texttt{0x31}), el acceso de seguridad (\texttt{0x27}) o la lectura de información de DTC (\texttt{0x19}), que permiten acceder a parámetros de configuración y calibración de las ECUs no expuestos por el estándar OBD-II.



\section{Reflexión sobre el uso de herramientas de inteligencia artificial}
\label{sec:uso_ia}

Durante el desarrollo de este trabajo se hizo uso de herramientas de inteligencia artificial generativa como asistente de programación, concretamente mediante el acceso a modelos de lenguaje de gran escala a través de interfaces de línea de comandos.

El uso de estas herramientas se circunscribió a tres contextos concretos: la generación de fragmentos de código repetitivos o bien definidos por un estándar (como las tablas de decodificación de PIDs OBD-II), la resolución de dudas puntuales sobre la API de las librerías utilizadas (\texttt{python-can}, \texttt{bless}, \texttt{react-native-ble-plx}) y la revisión de la sintaxis de LaTeX en secciones con tablas complejas.

En ningún caso se delegó en estas herramientas el diseño de la arquitectura del sistema, la implementación de los protocolos de comunicación ni la interpretación de los resultados experimentales.

La experiencia ha puesto de manifiesto tanto el valor como los límites de estas herramientas en el contexto de un proyecto de ingeniería. Son especialmente útiles para reducir el tiempo invertido en tareas de codificación mecánica y para obtener una primera orientación ante problemas con librerías poco documentadas. Sin embargo, la generación de código incorrecto sin diagnóstico claro del error sigue siendo habitual, y la verificación de cada sugerencia mediante prueba real es indispensable.